% GENERATED FILE. Edit canonical lessons, then run npm run generate:books.
% canonical-source-hash: fnv1a64:9aa27aa31849fba3
% canonical-lessons: HI-C412-vides, HI-C412-pardes, HI-C412-sarhad, HI-C412-utaarna, HI-C412-carhana

\chapter{Abroad and Away}
\label{ch:hi-abroad-and-away}
\hlchaptermodality{412}

\begin{chapteropening}
\textbf{By the end of this chapter:} \emph{I can talk about going abroad and living far from home, a frontier, and loading and unloading.}

Complete the last lesson of chapter 412: I can talk about going abroad and living far from home, a frontier, and loading and unloading.
\end{chapteropening}

\section[\texorpdfstring{\hi{विदेश} (videś)}{विदेश (videś)}]{\hi{विदेश} (videś) — abroad; a foreign country}

\label{lesson:HI-C412-vides}



\begin{quote}
Before the new one: say the Hindi for cancelled, then the Hindi for departure.
\end{quote}

\subsection*{You'll want to know: \hi{विदेश}}
\textbf{\hi{विदेश}} — \emph{videś} — \textquotedblleft{}abroad; a foreign country\textquotedblright{}.

\begin{quote}
\hi{मेरा} \hi{भाई} \hi{विदेश} \hi{में} \hi{काम} \hi{करता} \hi{है।}
\end{quote}

\begin{quote}
\emph{merā bhāī videś meṁ kām kartā hai.}
\end{quote}

\begin{quote}
My brother works abroad.
\end{quote}

Masculine.

\begin{cousinweb}
\textbf{\hi{विदेश}} (\emph{videś}) is a tatsama, Sanskrit \emph{videśa}, \textquotedblleft{}another country, abroad\textquotedblright{} (Monier-Williams), from \emph{vi-}, \textquotedblleft{}apart\textquotedblright{}, and \emph{deśa}, \textquotedblleft{}country\textquotedblright{}; \hi{विदेशी}, foreign, is its adjective.
\end{cousinweb}

\begin{grammarlens}[title={What you've built}]
One more word for getting around.
\end{grammarlens}

\subsection*{Your turn}
\begin{itemize}
\raggedright
  \item \emph{Say it:} \emph{videś}
  \item \emph{Say it:} \emph{videś}, once more
  \item \emph{Say it:} \emph{merā bhāī videś meṁ kām kartā hai.} — My brother works abroad.
  \item \emph{Recall:} say the Hindi for cancelled, then the Hindi for departure, then say \emph{videś} again
\end{itemize}

\begin{tcolorbox}[breakable,colback=teal!4,colframe=teal!35!black,title={Before you move on}]
What does \hi{विदेश} mean? (\textquotedblleft{}Abroad; a foreign country\textquotedblright{}.) Say it once more.
\end{tcolorbox}
\section[\texorpdfstring{\hi{परदेस} (pardes)}{परदेस (pardes)}]{\hi{परदेस} (pardes) — a land far from home}

\label{lesson:HI-C412-pardes}



\begin{quote}
Before the new one: say the Hindi for departure, then the Hindi for abroad.
\end{quote}

\subsection*{You'll want to know: \hi{परदेस}}
\textbf{\hi{परदेस}} — \emph{pardes} — \textquotedblleft{}a land far from home\textquotedblright{}.

\begin{quote}
\hi{परदेस} \hi{में} \hi{अपने} \hi{घर} \hi{की} \hi{बहुत} \hi{याद} \hi{आती} \hi{है।}
\end{quote}

\begin{quote}
\emph{pardes meṁ apne ghar kī bahut yād ātī hai.}
\end{quote}

\begin{quote}
Far from home, you miss your own house a lot.
\end{quote}

Masculine.

\begin{cousinweb}
\textbf{\hi{परदेस}} (\emph{pardes}) is a tadbhava of Sanskrit \emph{paradeśa}, \textquotedblleft{}another or a foreign country\textquotedblright{} (Monier-Williams). It may be another state of India: anywhere that is not home.
\end{cousinweb}

\begin{grammarlens}[title={What you've built}]
One more word for getting around.
\end{grammarlens}

\subsection*{Your turn}
\begin{itemize}
\raggedright
  \item \emph{Say it:} \emph{pardes}
  \item \emph{Say it:} \emph{pardes}, once more
  \item \emph{Say it:} \emph{pardes meṁ apne ghar kī bahut yād ātī hai.} — Far from home, you miss your own house a lot.
  \item \emph{Recall:} say the Hindi for departure, then the Hindi for abroad, then say \emph{pardes} again
\end{itemize}

\begin{tcolorbox}[breakable,colback=teal!4,colframe=teal!35!black,title={Before you move on}]
What does \hi{परदेस} mean? (\textquotedblleft{}A land far from home\textquotedblright{}.) Say it once more.
\end{tcolorbox}
\section[\texorpdfstring{\hi{सरहद} (sarhad)}{सरहद (sarhad)}]{\hi{सरहद} (sarhad) — a frontier; a border}

\label{lesson:HI-C412-sarhad}



\begin{quote}
Before the new one: say the Hindi for abroad, then the Hindi for a land far from home.
\end{quote}

\subsection*{You'll want to know: \hi{सरहद}}
\textbf{\hi{सरहद}} — \emph{sarhad} — \textquotedblleft{}a frontier; a border\textquotedblright{}.

\begin{quote}
\hi{सरहद} \hi{पर} \hi{बहुत} \hi{सिपाही} \hi{थे।}
\end{quote}

\begin{quote}
\emph{sarhad par bahut sipāhī the.}
\end{quote}

\begin{quote}
There were many soldiers at the frontier.
\end{quote}

Feminine.

\begin{cousinweb}
\textbf{\hi{सरहद}} (\emph{sarhad}) is usually traced to Persian \emph{sarhad}, \textquotedblleft{}a frontier\textquotedblright{}, literally the head of the boundary.
\end{cousinweb}

\begin{grammarlens}[title={What you've built}]
One more word for getting around.
\end{grammarlens}

\subsection*{Your turn}
\begin{itemize}
\raggedright
  \item \emph{Say it:} \emph{sarhad}
  \item \emph{Say it:} \emph{sarhad}, once more
  \item \emph{Say it:} \emph{sarhad par bahut sipāhī the.} — There were many soldiers at the frontier.
  \item \emph{Recall:} say the Hindi for abroad, then the Hindi for a land far from home, then say \emph{sarhad} again
\end{itemize}

\begin{tcolorbox}[breakable,colback=teal!4,colframe=teal!35!black,title={Before you move on}]
What does \hi{सरहद} mean? (\textquotedblleft{}A frontier; a border\textquotedblright{}.) Say it once more.
\end{tcolorbox}
\section[\texorpdfstring{\hi{उतारना} (utārnā)}{उतारना (utārnā)}]{\hi{उतारना} (utārnā) — to unload; to take down}

\label{lesson:HI-C412-utaarna}



\begin{quote}
Before the new one: say the Hindi for a land far from home, then the Hindi for a frontier.
\end{quote}

\subsection*{You'll want to know: \hi{उतारना}}
\textbf{\hi{उतारना}} — \emph{utārnā} — \textquotedblleft{}to unload; to take down\textquotedblright{}.

\begin{quote}
\hi{कुली} \hi{ने} \hi{ट्रेन} \hi{से} \hi{हमारा} \hi{सामान} \hi{उतारा।}
\end{quote}

\begin{quote}
\emph{kulī ne ṭren se hamārā sāmān utārā.}
\end{quote}

\begin{quote}
The porter unloaded our luggage from the train.
\end{quote}

Past, with \hi{ने} (he or she …): \hi{उसने} \hi{छत} \hi{से} \hi{कपड़े} \hi{उतारे}, \emph{usne chat se kapṛe utāre}.

\begin{cousinweb}
\textbf{\hi{उतारना}} (\emph{utārnā}) is the causative partner of \hi{उतरना}, from Sanskrit \emph{uttārayati}, \textquotedblleft{}makes descend\textquotedblright{} (CDIAL 1770), which lists Hindi \emph{utārnā}, \textquotedblleft{}to make descend\textquotedblright{}.
\end{cousinweb}

\begin{grammarlens}[title={What you've built}]
One more word for getting around.
\end{grammarlens}

\subsection*{Your turn}
\begin{itemize}
\raggedright
  \item \emph{Say it:} \emph{utārnā}
  \item \emph{Say it:} \emph{utārnā}, once more
  \item \emph{Say it:} \emph{kulī ne ṭren se hamārā sāmān utārā.} — The porter unloaded our luggage from the train.
  \item \emph{Recall:} say the Hindi for a land far from home, then the Hindi for a frontier, then say \emph{utārnā} again
\end{itemize}

\begin{tcolorbox}[breakable,colback=teal!4,colframe=teal!35!black,title={Before you move on}]
What does \hi{उतारना} mean? (\textquotedblleft{}To unload; to take down\textquotedblright{}.) Say it once more.
\end{tcolorbox}
\section[\texorpdfstring{\hi{चढ़ाना} (caṛhānā)}{चढ़ाना (caṛhānā)}]{\hi{चढ़ाना} (caṛhānā) — to put on board; to put up}

\label{lesson:HI-C412-carhana}



\begin{quote}
Before the new one: say the Hindi for a frontier, then the Hindi for to unload.
\end{quote}

\subsection*{You'll want to know: \hi{चढ़ाना}}
\textbf{\hi{चढ़ाना}} — \emph{caṛhānā} — \textquotedblleft{}to put on board; to put up\textquotedblright{}.

\begin{quote}
\hi{हमने} \hi{अटैची} \hi{बस} \hi{की} \hi{छत} \hi{पर} \hi{चढ़ाई।}
\end{quote}

\begin{quote}
\emph{hamne aṭaicī bas kī chat par caṛhāī.}
\end{quote}

\begin{quote}
We put the suitcase up on the roof of the bus.
\end{quote}

Past, with \hi{ने} (he or she …): \hi{उसने} \hi{सामान} \hi{ट्रेन} \hi{में} \hi{चढ़ाया}, \emph{usne sāmān ṭren meṁ caṛhāyā}.

\begin{cousinweb}
\textbf{\hi{चढ़ाना}} (\emph{caṛhānā}) is the causative partner of \hi{चढ़ना}, from \emph{*cāḍhayati}, \textquotedblleft{}raises\textquotedblright{}, a form CDIAL reconstructs (CDIAL 4578), which lists Hindi \emph{caṛhānā}.
\end{cousinweb}

\begin{grammarlens}[title={What you've built}]
One more word for getting around.
\end{grammarlens}

\subsection*{Your turn}
\begin{itemize}
\raggedright
  \item \emph{Say it:} \emph{caṛhānā}
  \item \emph{Say it:} \emph{caṛhānā}, once more
  \item \emph{Say it:} \emph{hamne aṭaicī bas kī chat par caṛhāī.} — We put the suitcase up on the roof of the bus.
  \item \emph{Recall:} say the Hindi for a frontier, then the Hindi for to unload, then say \emph{caṛhānā} again
\end{itemize}

\begin{tcolorbox}[breakable,colback=teal!4,colframe=teal!35!black,title={Before you move on}]
What does \hi{चढ़ाना} mean? (\textquotedblleft{}To put on board; to put up\textquotedblright{}.) Say it once more.
\end{tcolorbox}
